% Source: Paper_A_新版完整论文提纲_vFinal_Draft.md §13
\section{试验分布与生产标准化}

\subsection{分层结果}

分别报告：
\[
V_{\pi,o},\qquad V_{\pi,s}.
\]

\subsection{50:50 设计平均}

\[
V_\pi^{\mathrm{design}}
=
0.5V_{\pi,o}+0.5V_{\pi,s}.
\]

只能解释为 balanced experimental mixture。

\subsection{生产标准化}

若目标生产池比例为 \(p_o,p_s\)：
\[
V_\pi^{\mathrm{prod}}
=
p_oV_{\pi,o}+p_sV_{\pi,s}.
\]

比例来自独立自然任务池，不从 50:50 试验样本估计。

上式只直接适用于 delivery-adjusted quality、non-delivery/severe rate 和平均成本等
无条件指标。resolved-only quality 必须按分子与分母标准化：

\[
Q_{\pi,\mathrm{resolved}}^{prod}
=
\frac{p_o r_o q_o+p_s r_s q_s}{p_o r_o+p_s r_s},
\]

其中 \(r_o,r_s\) 为各层 resolved rate，不能机械加权两个条件性均值。

若没有唯一比例，报告：
\begin{itemize}
  \item \texttt{80:20}
  \item \texttt{60:40}
  \item \texttt{50:50}
  \item \texttt{30:70}
\end{itemize}
情景分析。
